\documentclass[10pt,a4paper]{amsart}
\usepackage[margin=19mm]{geometry}
\usepackage{amsmath,amssymb,mathtools}
\usepackage[scr=rsfs,cal=euler]{mathalfa}
\usepackage{xfrac}
\usepackage[unicode,hidelinks,bookmarks=false]{hyperref}
\newcommand{\ie}{i.e.\ }
\newcommand{\cf}{cf.\ }
\newcommand{\resp}{respectively\ }
\newcommand{\Subsection}{\subsection}
\providecommand{\nobreakdash}{\nobreak\hbox{-}}
\setlength{\emergencystretch}{2em}
\multlinegap0pt
\usepackage{amssymb}
\usepackage{dsfont}
\newtheorem{theorem}{Theorem}
\newcommand{\GC}{\mathcal{G}}
\newcommand{\IC}{\mathcal{I}}
\newcommand{\KC}{\mathcal{K}}
\newcommand{\id}{\mathrm{id}}
\newcommand{\unit}{\mathds{1}}
\title{The Small-Gain Condition for Infinite Networks Modeled on $\ell^{\infty}$-Spaces}
\author{Christoph Kawan\, \, }
\date{Source: arXiv:2503.03925v7 (CC BY 4.0)}
\begin{document}
\maketitle

\noindent\emph{Source and licence.} The Small-Gain Condition for Infinite Networks Modeled on $\ell^{\infty}$-Spaces. Version arXiv:2503.03925v7 (CC BY 4.0), distributed by arXiv under the Creative Commons Attribution 4.0 International licence, http://creativecommons.org/licenses/by/4.0/, as recorded in the arXiv licence metadata for this exact version and shown on the abstract page https://arxiv.org/abs/2503.03925v7. Full licence notice: \texttt{licenses/arxiv-2503.03925-CC-BY-4.0.txt}.\par\smallskip

\noindent\textit{Editorial scope.} Counted unit: the article's theorem thm\_finite\_general\_case (source0.tex lines 1292--1300). The article's complete original proof of it is delivered with it (source0.tex lines 1302--1345, one \texttt{proof} environment of 44 lines). No mathematics is written, repaired or re-derived: the statement and the proof are the article's own lines, in the article's own notation, and no retained line is substituted or re-flowed.  Retained uncounted as the source-local context the proof needs: lines 1016--1016; lines 1277--1279.  Nothing was omitted from any retained span, so the package carries no omission record.  No retained line contains a citation, so the wrapper carries no bibliography and no bibliography entry.  No retained line contains a figure environment, an image-inclusion command or drawing code, so no image is delivered and the package declares no asset dependency.  Editorial formatting only, in addition: an AMS \texttt{amsart} preamble that restates the article's own declarations and macros used by the retained lines, namely the environments \texttt{theorem} on separate counters, together with the standard packages those lines need.  The retained environments print in this document as Theorem 1 in order of appearance, whereas the article prints the same items with the numbering of its own full document.  The complete native source of the article as distributed in the arXiv e-print for this exact version is bundled unchanged as \texttt{sources/174293/source0.tex}.
\section{Main results}\label{sec_final_theorem}

\begin{equation}\label{eq_mafs_uniform_cont}
  \mu_i(s^2) - \mu_i(s^1) \leq \gamma(\|s^2 - s^1\|) \mbox{\quad for all\ } i \in \IC.
\end{equation}

\begin{theorem}\label{thm_finite_general_case}
Assume that $|\IC| < \infty$ and consider a gain operator with uniformly continuous MAFs. Then the following statements are equivalent:%
\begin{enumerate}
\item[(a)] There exists a path of strict decay for $\Gamma$.%
\item[(b)] There exists $\rho \in \KC_{\infty}$ such that $\Sigma(\Gamma_{\rho})$ is UGAS.
\item[(c)] There exists $\rho \in \KC_{\infty}$ such that $\Gamma_{\rho}$ satisfies the $\oplus$-MBI property.%
\item[(d)] There exists $\rho \in \KC_{\infty}$ such that $\Gamma_{\rho}$ satisfies the NJI condition.%
\end{enumerate}
\end{theorem}

\begin{proof}
Looking at our general results in Section \ref{sec_final_theorem} and observing that (b) trivially implies (d), we can easily see that the proof is completed by showing that (d) implies (c). To this end, consider $s \in \ell^{\infty}_+(\IC)$ and $r \geq 0$ with $s \leq r\unit \oplus \Gamma(s)$. Without loss of generality, we can assume that $s > 0$. We split the index set $\IC$ into the two subsets%
\begin{equation*}
  \IC^1 := \left\{ i \in \IC : s_i \leq r \right\},\quad \IC^2 := \IC \backslash \IC^1.%
\end{equation*}
By assumption, there exists at least one index $i$ with $\Gamma_i(s) < s_i$, which implies $s_i \leq r$. Hence, $\IC^1$ is nonempty. We define the $\KC_{\infty}$-function%
\begin{equation*}
  \varphi_1 := (\rho \circ \xi \circ \eta)^{-1} \circ \gamma \circ \left(\max_{ji \in E(\GC)}\gamma_{ij}\right),%
\end{equation*}
where $\rho$ comes from (a), $\xi$ comes from (M1), $\eta \leq \gamma_{ij}$ for all $ji \in E(\GC)$ with $\eta \in \KC_{\infty}$ (such $\eta$ trivially exists), and $\gamma$ comes from \eqref{eq_mafs_uniform_cont}. Without loss of generality, we can assume that%
\begin{equation}\label{eq_finite_general_wlogs}
  \varphi_1(r) > r \mbox{\quad and \quad} (\rho \circ \xi \circ \eta)^{-1}(r) > r \mbox{\quad for all\ } r > 0.%
\end{equation}
Let us assume to the contrary that $s_{|\IC^2} \geq \varphi_1(r)\unit$. We want to show that this implies%
\begin{equation*}
  (\id + \rho) \circ \Gamma_i(s_{|\IC^2}) \geq \Gamma_i(s) \mbox{\quad for all\ } i \in \IC^2,% 
\end{equation*}
which in turn implies $\Gamma_{\rho}(s_{|\IC^2}) \geq s_{|\IC^2}$ in contradiction to (a). We fix $i \in \IC^2$ and define%
\begin{equation*}
  a_i := [\gamma_{ij}(s_j)]_{j \in \IC_i \cap \IC^2}.%
\end{equation*}
We can assume without loss of generality that $a_i \neq 0$. Indeed, $a_i = 0$ is only possible if $\IC_i \cap \IC^2 = \emptyset$, implying%
\begin{align*}
  s_i &\leq \max\{r,\mu_i([\gamma_{ij}(r)]_{j \in \IC_i})\} \leq \max\left\{r, \gamma \circ  \left(\max_{ji \in E(\GC)}\gamma_{ij}\right)(r) \right\} \stackrel{\eqref{eq_finite_general_wlogs}}{<} \varphi_1(r)%			
\end{align*}
in contradiction to our assumption hat $s_{|\IC^2} \geq \varphi_1(r)\unit$. The inequality to be shown can then be written as%
\begin{equation*}
  \rho(\mu_i(a_i)) \geq \mu_i([\gamma_{ij}(r)]_{j \in \IC_i \cap \IC^1} + a_i) - \mu_i(a_i).%
\end{equation*}
Using \eqref{eq_mafs_uniform_cont}, we see that a sufficient condition for this inequality to hold is%
\begin{equation*}
  \rho(\mu_i(a_i)) \geq \gamma( \|[\gamma_{ij}(r)]_{j \in \IC_i \cap \IC^1}\| ).%
\end{equation*}
This last inequality follows from the definition of $\varphi_1$, since%
\begin{align*}
  \rho(\mu_i(a_i)) &\geq \rho(\xi(\|a_i\|)) \geq \rho \circ \xi \circ \eta \circ \varphi_1(r) \\
	                 &= \gamma \circ \left(\max_{ji \in E(\GC)}\gamma_{ij}\right)(r) \geq \gamma( \|[\gamma_{ij}(r)]_{j \in \IC_i \cap \IC^1}\| ).%
\end{align*}
Hence, there exists at least one index $i \in \IC^2$ such that $s_i  < \varphi_1(r)$. We can split the index set $\IC$ again into%
\begin{equation*}
  \IC^1 := \left\{ i \in \IC : s_i \leq \max\{r,\varphi_1(r)\} \right\},\quad \IC^2 := \IC \backslash \IC^1,%
\end{equation*}
observing that the new $\IC^1$ contains at least one more element than the old one. With the same arguments as before, we can define another $\KC_{\infty}$-function $\varphi_2$ such that $s_i < \varphi_2(r)$ must hold for at least one $i \in \IC^2$. In this way, we can continue inductively until $\IC^2 = \emptyset$ (using that there are only finitely many indices). We thus find a $\KC_{\infty}$-function $\varphi$ such that $\|s\| \leq \varphi(r)$. Note that we have proved that $\Gamma$ satisfies the $\oplus$-MBI property (and not $\Gamma_{\rho}$). By choosing a uniformly continuous $\KC_{\infty}$-function $\rho' < \rho$, the proof also works for $\Gamma_{\rho'}$ in place of $\Gamma$.
\end{proof}

\end{document}
